\documentclass[10pt,a4paper]{article}

\usepackage[T1]{fontenc}
\usepackage[utf8]{inputenc}
\usepackage{lmodern}
\usepackage{microtype}
\usepackage[a4paper,margin=18mm,headheight=21pt,footskip=12mm]{geometry}
\usepackage{xcolor}
\usepackage{amsmath,amssymb}
\usepackage{booktabs,tabularx,array}
\usepackage{enumitem}
\usepackage{fancyhdr}
\usepackage{tcolorbox}
\usepackage{needspace}
\usepackage{hyperref}

\definecolor{navy}{HTML}{10263F}
\definecolor{blue}{HTML}{1E6A8D}
\definecolor{teal}{HTML}{1B8A89}
\definecolor{sky}{HTML}{EAF4F8}
\definecolor{mist}{HTML}{F4F7F9}
\definecolor{line}{HTML}{CBD6DC}
\definecolor{ink}{HTML}{16222A}
\definecolor{muted}{HTML}{536873}

\hypersetup{colorlinks=true,linkcolor=blue,urlcolor=blue,pdfauthor={Independent practice material},pdftitle={QRT-style Quant Research Mock Assessment}}
\color{ink}
\setlength{\parindent}{0pt}
\setlength{\parskip}{4pt}
\setlist[enumerate]{label=\textbf{\Alph*.},leftmargin=8mm,itemsep=1.5pt,topsep=3pt,parsep=0pt}

\pagestyle{fancy}
\fancyhf{}
\lhead{\textcolor{muted}{QRT-style Quant Research Mock}}
\rhead{\textcolor{muted}{82 minutes}}
\cfoot{\textcolor{muted}{\thepage}}
\renewcommand{\headrulewidth}{0.4pt}
\renewcommand{\headrule}{\hbox to\headwidth{\color{line}\leaders\hrule height \headrulewidth\hfill}}

\newcommand{\sectionbanner}[2]{%
  \vspace{3mm}
  {\color{navy}\Large\bfseries #1}\hfill
  {\color{blue}\bfseries #2}
  \par\vspace{1mm}\color{line}\rule{\linewidth}{0.8pt}\color{ink}\vspace{2mm}
}

\newtcolorbox{questionbox}[2]{
  colback=white,
  colframe=line,
  colbacktitle=mist,
  coltitle=navy,
  fonttitle=\bfseries,
  title={Question #1\hfill\textcolor{blue}{\small #2}},
  boxrule=0.7pt,
  arc=1.5mm,
  left=3mm,right=3mm,top=2.2mm,bottom=2.2mm,
  before skip=7pt,after skip=7pt
}

\newtcolorbox{notebox}{
  colback=sky,
  colframe=teal,
  boxrule=0.7pt,
  arc=1.5mm,
  left=3mm,right=3mm,top=2.5mm,bottom=2.5mm
}

\begin{document}

\thispagestyle{empty}
\vspace*{10mm}
{\color{blue}\rule{18mm}{2.2pt}}\\[7mm]
{\fontsize{28}{32}\selectfont\bfseries\color{navy} Quant Research\\Mock Assessment}\\[4mm]
{\Large\color{blue} QRT-style practice edition}\par

\vspace{10mm}
\begin{tcolorbox}[colback=navy,colframe=navy,arc=2mm,left=5mm,right=5mm,top=4mm,bottom=4mm]
  \color{white}
  \begin{tabularx}{\linewidth}{>{\bfseries}lX>{\bfseries}r}
    Total time & 82 minutes & 24 items \\
    Coding & 1 data-science question in the companion notebook & 42 min suggested \\
    Theory & 10 math + 5 ML fundamentals + 8 ML advanced & 40 min suggested
  \end{tabularx}
\end{tcolorbox}

\vspace{9mm}
\begin{tabularx}{\linewidth}{@{}>{\color{muted}\bfseries}l X@{}}
Candidate & \rule{0.78\linewidth}{0.4pt} \\[6mm]
Date & \rule{0.78\linewidth}{0.4pt} \\[6mm]
Start time & \rule{0.34\linewidth}{0.4pt}\hspace{9mm}\textcolor{muted}{End time}\hspace{3mm}\rule{0.27\linewidth}{0.4pt}
\end{tabularx}

\vfill
\begin{notebox}
\textbf{Independent practice material.} This is an original mock, not an official or leaked Qube Research \& Technologies assessment. Public candidate reports were used only to calibrate section structure and topic coverage. The invitation and on-screen rules for your actual assessment take priority.
\end{notebox}

\vspace{5mm}
{\small\color{muted}Package: \texttt{QRT\_mock\_test} \hfill Prepared 28 October 2026}
\newpage

\sectionbanner{Instructions}{Read before starting}

\begin{enumerate}[label=\arabic*.,leftmargin=7mm,itemsep=3pt]
  \item Start one 82-minute timer. A suggested allocation is 42 minutes for the notebook, 18 minutes for Questions 1--10, 8 minutes for Questions 11--15, and 14 minutes for Questions 16--23.
  \item The companion notebook contains the single coding question. This PDF contains the other 23 questions.
  \item For \textbf{single-answer} questions, choose exactly one option. For \textbf{select-all} questions, choose every correct option and no incorrect option. Suggested scoring is one point only for an exact match.
  \item Unless a question says otherwise, use ordinary exact arithmetic and standard statistical conventions. A calculator is optional for this mock.
  \item Do not open the \texttt{solutions/} folder until the timer ends. There is no penalty for guessing in this practice version.
\end{enumerate}

\vspace{4mm}
\begin{notebox}
\textbf{Format confidence.} The 82-minute duration and 1/10/5/8 counts come from the user-provided brief. A recent third-party synthesis independently reports the same four section counts, but reports 90 minutes and is not an official QRT publication. Public candidate accounts support Python, statistics, OLS, regularization, and tree-ensemble coverage while also showing role and office variation. Platform, proctoring, scoring, navigation, and calculator rules remain unverified and are not simulated here.
\end{notebox}

\vspace{5mm}
\textbf{Public references used for calibration}

\begin{itemize}[leftmargin=6mm,itemsep=2pt]
  \item \href{https://www.qube-rt.com/careers/research-and-trading/london/students-and-new-graduates/}{QRT official Students and New Grads page} -- official careers context, no assessment blueprint.
  \item \href{https://www.tradermath.org/articles/qube-research-technologies-interview-guide}{TraderMath QRT interview guide} -- third-party candidate-report synthesis; same section counts, different timing.
  \item \href{https://www.glassdoor.com/Interview/Qube-Research-and-Technologies-Interview-Questions-E3226090.htm}{Glassdoor QRT interview reports} and \href{https://www.wallstreetoasis.com/company/qube-research-and-technologies/interview}{Wall Street Oasis QRT entries} -- anonymous, variable experiences.
\end{itemize}

\vspace{5mm}
\textbf{Answer sheet}

\renewcommand{\arraystretch}{1.55}
\begin{center}
\begin{tabular}{c p{24mm} c p{24mm} c p{24mm}}
\toprule
\textbf{Q} & \textbf{Answer} & \textbf{Q} & \textbf{Answer} & \textbf{Q} & \textbf{Answer} \\
\midrule
1 & \rule{20mm}{0.3pt} & 9 & \rule{20mm}{0.3pt} & 17 & \rule{20mm}{0.3pt} \\
2 & \rule{20mm}{0.3pt} & 10 & \rule{20mm}{0.3pt} & 18 & \rule{20mm}{0.3pt} \\
3 & \rule{20mm}{0.3pt} & 11 & \rule{20mm}{0.3pt} & 19 & \rule{20mm}{0.3pt} \\
4 & \rule{20mm}{0.3pt} & 12 & \rule{20mm}{0.3pt} & 20 & \rule{20mm}{0.3pt} \\
5 & \rule{20mm}{0.3pt} & 13 & \rule{20mm}{0.3pt} & 21 & \rule{20mm}{0.3pt} \\
6 & \rule{20mm}{0.3pt} & 14 & \rule{20mm}{0.3pt} & 22 & \rule{20mm}{0.3pt} \\
7 & \rule{20mm}{0.3pt} & 15 & \rule{20mm}{0.3pt} & 23 & \rule{20mm}{0.3pt} \\
8 & \rule{20mm}{0.3pt} & 16 & \rule{20mm}{0.3pt} & & \\
\bottomrule
\end{tabular}
\end{center}

\newpage
\sectionbanner{Section A: Mathematics and Statistics}{Questions 1--10 | 18 min suggested}

\Needspace{9\baselineskip}
\begin{questionbox}{1}{Single answer}
A market is in a high-volatility regime on 20\% of days. A detector flags 60\% of high-volatility days and 10\% of all other days. Given that the detector flags today, what is the probability that today is a high-volatility day?
\begin{enumerate}
  \item 0.20
  \item 0.40
  \item 0.60
  \item 0.75
\end{enumerate}
\end{questionbox}

\Needspace{10\baselineskip}
\begin{questionbox}{2}{Single answer}
Daily event counts are conditionally Poisson with unknown rate $\lambda$. The prior is $\lambda\sim\mathrm{Gamma}(3,2)$, using shape and rate. Over three equal observation periods the counts are 4, 2, and 5. What is the posterior mean of $\lambda$?
\begin{enumerate}
  \item $11/3$
  \item $7/3$
  \item $14/5$
  \item $16/5$
\end{enumerate}
\end{questionbox}

\Needspace{9\baselineskip}
\begin{questionbox}{3}{Single answer}
Suppose $X_1,\ldots,X_n$ are i.i.d. $\mathrm{Uniform}(0,\theta)$ and the sample mean is 4.2. What is the method-of-moments estimate of $\theta$?
\begin{enumerate}
  \item 2.1
  \item 4.2
  \item 6.3
  \item 8.4
\end{enumerate}
\end{questionbox}

\Needspace{9\baselineskip}
\begin{questionbox}{4}{Single answer}
An OLS regression with an intercept has fitted slope $-0.8$. The sample means are $\bar{x}=1.5$ and $\bar{y}=0.7$. What is the fitted intercept?
\begin{enumerate}
  \item $-1.9$
  \item $-0.5$
  \item 0.1
  \item 1.9
\end{enumerate}
\end{questionbox}

\Needspace{9\baselineskip}
\begin{questionbox}{5}{Single answer}
A score is normally distributed with mean 50 and standard deviation 4. What is the z-score of an observation equal to 58?
\begin{enumerate}
  \item 1
  \item 2
  \item 4
  \item 8
\end{enumerate}
\end{questionbox}

\Needspace{10\baselineskip}
\begin{questionbox}{6}{Single answer}
What is the rank of the matrix
\[
A=\begin{pmatrix}
1&2&0\\
0&1&1\\
1&3&1
\end{pmatrix}?
\]
\begin{enumerate}
  \item 0
  \item 1
  \item 2
  \item 3
\end{enumerate}
\end{questionbox}

\Needspace{10\baselineskip}
\begin{questionbox}{7}{Single answer}
A joint density is $f(x,y)=c(x+2y)$ for $0\leq x\leq1$ and $0\leq y\leq1$, and zero elsewhere. What value of $c$ normalizes the density?
\begin{enumerate}
  \item $1/2$
  \item $2/3$
  \item 1
  \item $3/2$
\end{enumerate}
\end{questionbox}

\Needspace{10\baselineskip}
\begin{questionbox}{8}{Single answer}
A chi-squared goodness-of-fit test uses seven categories. Two distribution parameters are estimated from the same data. Assuming regularity conditions, how many degrees of freedom does the reference chi-squared distribution have?
\begin{enumerate}
  \item 2
  \item 3
  \item 4
  \item 6
\end{enumerate}
\end{questionbox}

\Needspace{10\baselineskip}
\begin{questionbox}{9}{Single answer}
Let $\operatorname{Var}(X)=4$, $\operatorname{Var}(Y)=9$, and $\operatorname{Cov}(X,Y)=3$. What is $\operatorname{Var}(2X-Y)$?
\begin{enumerate}
  \item 7
  \item 11
  \item 13
  \item 37
\end{enumerate}
\end{questionbox}

\Needspace{13\baselineskip}
\begin{questionbox}{10}{Select all that apply}
In the linear model $y=X\beta+\varepsilon$, which conditions are required for the usual OLS coefficient vector to be uniquely defined and conditionally unbiased given $X$? Assume the relevant moments exist.
\begin{enumerate}
  \item $X$ has full column rank.
  \item $\mathbb{E}[\varepsilon\mid X]=0$.
  \item $\operatorname{Var}(\varepsilon\mid X)=\sigma^2I$.
  \item $\varepsilon\mid X$ is normally distributed.
\end{enumerate}
\end{questionbox}

\newpage
\sectionbanner{Section B: Machine Learning Fundamentals}{Questions 11--15 | 8 min suggested}

\Needspace{11\baselineskip}
\begin{questionbox}{11}{Single answer}
Which workflow most clearly leaks validation information into model training?
\begin{enumerate}
  \item Split first, fit a scaler on the training set, and transform both sets with it.
  \item Tune a model using cross-validation inside the training set, then evaluate once on the test set.
  \item Standardize all rows first, then split the standardized data into training and validation sets.
  \item Fit a feature encoder on each training fold and apply it to that fold's validation rows.
\end{enumerate}
\end{questionbox}

\Needspace{14\baselineskip}
\begin{questionbox}{12}{Select all that apply}
Which statements about linear-model regularization are generally correct?
\begin{enumerate}
  \item Increasing the penalty strength typically increases bias.
  \item Ridge regression commonly sets many coefficients exactly to zero.
  \item Lasso can set some coefficients exactly to zero.
  \item Numeric features should usually be put on comparable scales before applying either penalty.
  \item Increasing the penalty strength typically increases estimator variance.
\end{enumerate}
\end{questionbox}

\Needspace{11\baselineskip}
\begin{questionbox}{13}{Single answer}
Only 0.5\% of observations are positive, and finding positives is the primary objective. Which metric is usually more informative than raw accuracy for comparing ranking models?
\begin{enumerate}
  \item Mean squared error
  \item Area under the precision-recall curve
  \item Adjusted $R^2$
  \item Explained variance
\end{enumerate}
\end{questionbox}

\Needspace{11\baselineskip}
\begin{questionbox}{14}{Single answer}
You must validate a model that predicts next month's outcome from historical time-series features. Which split best respects the information set available at prediction time?
\begin{enumerate}
  \item Randomly shuffle all rows before five-fold cross-validation.
  \item Train on later dates and validate on earlier dates.
  \item Use expanding training windows followed by later validation blocks.
  \item Duplicate rare months before making a random split.
\end{enumerate}
\end{questionbox}

\Needspace{10\baselineskip}
\begin{questionbox}{15}{Single answer}
Which method is unsupervised?
\begin{enumerate}
  \item Logistic regression
  \item K-means clustering
  \item Gradient-boosted classification
  \item Quantile regression
\end{enumerate}
\end{questionbox}

\newpage
\sectionbanner{Section C: Machine Learning Advanced}{Questions 16--23 | 14 min suggested}

\Needspace{12\baselineskip}
\begin{questionbox}{16}{Select all that apply}
Which mechanisms directly help decorrelate the trees in a standard random forest?
\begin{enumerate}
  \item Training trees on bootstrap samples.
  \item Considering a random subset of features at each split.
  \item Forcing every tree to use the same best feature at its root.
  \item Replacing the ensemble average with one deepest tree.
\end{enumerate}
\end{questionbox}

\Needspace{11\baselineskip}
\begin{questionbox}{17}{Single answer}
In gradient boosting with squared-error loss, what target does each new weak learner approximately fit?
\begin{enumerate}
  \item The original response with no reference to the current ensemble.
  \item The current residuals, equivalently the negative loss gradients.
  \item The largest principal component of the feature matrix.
  \item The predictions of a separately trained random forest.
\end{enumerate}
\end{questionbox}

\Needspace{13\baselineskip}
\begin{questionbox}{18}{Select all that apply}
Which statements about the Adam optimizer are correct?
\begin{enumerate}
  \item It maintains exponential moving averages of first and second gradient moments.
  \item It guarantees convergence to a global minimum for a non-convex neural network.
  \item It creates parameter-specific adaptive step scaling.
  \item It removes the need to choose a learning rate.
  \item It cannot be used with mini-batches.
\end{enumerate}
\end{questionbox}

\Needspace{12\baselineskip}
\begin{questionbox}{19}{Single answer}
You use PCA before a classifier and estimate performance by cross-validation. Which implementation avoids leakage?
\begin{enumerate}
  \item Fit PCA once on all data, then cross-validate the classifier.
  \item Fit PCA once on all features but exclude the target column.
  \item Within each fold, fit PCA on the training portion and transform the validation portion with that fitted PCA.
  \item Fit PCA independently on the validation portion so it has its own components.
\end{enumerate}
\end{questionbox}

\Needspace{13\baselineskip}
\begin{questionbox}{20}{Select all that apply}
Two predictors are strongly correlated and contain similar signal. What can happen to ordinary single-feature permutation importance?
\begin{enumerate}
  \item Each feature's individual importance may be understated because the other can substitute for it.
  \item The importance of both features must always be overstated.
  \item Permutation importance automatically removes the correlation problem.
  \item Permuting the correlated features as a group is a useful sensitivity check.
\end{enumerate}
\end{questionbox}

\Needspace{11\baselineskip}
\begin{questionbox}{21}{Single answer}
What is the main purpose of nested cross-validation when many hyperparameters are compared?
\begin{enumerate}
  \item Make the final model sparse by construction.
  \item Use inner folds for selection and outer folds for a less biased estimate of the full tuning procedure.
  \item Guarantee that the chosen hyperparameters are globally optimal.
  \item Remove the need for a time-aware split in time-series data.
\end{enumerate}
\end{questionbox}

\Needspace{11\baselineskip}
\begin{questionbox}{22}{Single answer}
A binary classifier's scores are passed through a strictly increasing nonlinear transformation. Which statement is correct?
\begin{enumerate}
  \item ROC-AUC is unchanged, but calibration and log loss may change.
  \item ROC-AUC and log loss are both necessarily unchanged.
  \item ROC-AUC necessarily decreases.
  \item Calibration is necessarily unchanged because the ranking is unchanged.
\end{enumerate}
\end{questionbox}

\Needspace{11\baselineskip}
\begin{questionbox}{23}{Single answer}
After deployment, the marginal distribution of features remains approximately stable, but the relationship between features and the target changes. What is this most directly called?
\begin{enumerate}
  \item Covariate shift only
  \item Label shift only
  \item Concept drift
  \item Bootstrap instability
\end{enumerate}
\end{questionbox}

\vfill
\begin{notebox}
\textbf{Stop.} Record your answers, save the notebook, and end the timer before opening \texttt{solutions/answer\_key.md}.
\end{notebox}

\end{document}
